\documentclass[11pt]{article}
\usepackage{fontspec}
\usepackage[french]{babel}
\usepackage{amsmath}
\usepackage{amssymb}
\usepackage{amsthm}
\usepackage{mathrsfs}
\usepackage{tikz-cd}
\usepackage{enumitem}

\newcommand{\printedpage}[1]{\marginpar{\footnotesize p.~#1}}
\newcommand{\msp}[1]{\marginpar{\footnotesize [#1]}}
\newcommand{\uHom}{\underline{\mathrm{Hom}}}
\newcommand{\Ens}{\mathrm{Ens}}
\newcommand{\Fais}{\mathrm{Fais}}
\newcommand{\pstop}{\mathrm{pstop}}

\begin{document}

\printedpage{} % no page number printed
% head: (none)

\begin{center}
{\Large\textsf{LES DÉRIVATEURS}}

\bigskip

\textbf{ALEXANDRE GROTHENDIECK}

\bigskip\bigskip

\textbf{Chapitre XVIII}

\medskip

\textbf{Catégories et ensembles accessibles}
\end{center}

\newpage

\printedpage{} % no page number printed
% head: (none)

Ce texte a été déchiffré et transcrit en \LaTeXe{} par M. Künzer. Il a été édité par M. Künzer, J. Malgoire, et G. Maltsiniotis. La transcription est aussi fidèle que possible au manuscrit. Seules quelques corrections mineures évidentes ont été effectuées. Pour les rares ajouts ou commentaires des éditeurs, ainsi que pour la numérotation originale des pages du manuscrit, les caractères de machine à écrire \texttt{[typewriter]} entre crochets sont utilisés. Quelques notes de bas de page ajoutées par les éditeurs sont signalées par «~N.~Éd.~». La numérotation des pages de ce chapitre commence à la page 50, car à l'origine ce chapitre était la suite du chapitre XIX qui se termine au milieu de la page 50.

Cette édition est provisoire. Les remarques, commentaires et corrections sont bienvenus. Envoyer un message à :

\textbf{georges.maltsiniotis@imj-prg.fr}

\hfill G. Maltsiniotis\hspace*{2cm}

\newpage

\printedpage{1}
% head: A. Grothendieck : Dérivateurs, Ch. XVIII, Édition des Universités de Montpellier II et Paris VII
% foot: Version 17 décembre 2025

\msp{page 50}

\section*{4\quad Paratopos et l'opération $\mathcal{M}\boxtimes\mathcal{N}$}

\subsection*{4.1\quad Pseudo-topos. Problème du produit.}

J'appelle \textit{catégorie pseudo-$\mathfrak{U}$-topos}, ou simplement \textit{pseudo-topos}, une $\mathfrak{U}$-catégorie $\mathcal{M}$ stable par $\mathfrak{U}$-limites inductives, et qui admet une \textit{petite} sous-catégorie $\mathcal{M}_0$ \textit{strictement génératrice} (i.e. génératrice par épimorphismes stricts). Alors $\mathcal{M}$ admet aussi des $\mathfrak{U}$-limites projectives (SGA 4 I 8.12.7). (Mais l'exemple de (Gr), cf. loc. cit. 8.12.9, montre que $\mathcal{M}^\circ$ n'est par forcément, pour autant, un pseudo-$\mathfrak{U}$-topos, faute à $\mathcal{M}$ d'admettre une petite sous-catégorie strictement \textit{co}génératrice. Mais le théorème 2.6 \texttt{[du chapitre XIX]} donne des conditions sur $\mathcal{M}$ (surtout du type exactitude) % [sic] par forcément -> pas forcément, misprint; also the parenthesis opened at "(Mais" is never closed
\msp{page 51}
pour que $\mathcal{M}$ admette assez d'injectifs -- ce qui impliquera alors l'existence d'une petite sous-catégorie strictement cogénératrice, donc que $\mathcal{M}^\circ$ est également un pseudo-topos (\footnote{C'est le cas notamment si $\mathcal{M}$ est un topos.}).

Je visualise les catégories pseudo-topos comme étant la catégorie des “faisceaux” sur un objet $\mathcal{X}$, visualisé comme une sorte d'“espace” que je distingue de la catégorie $\mathcal{M}=\Fais(\mathcal{X})$ des faisceaux sur $\mathcal{X}$. C'est $\mathcal{X}$ qui sera pour ceci le “pseudo-topos”. Les pseudo-topos forment une 2-catégorie, la variance dans $\mathcal{X}$ étant opposée de celle dans $\Fais(\mathcal{X})=\mathcal{M}$. De façon précise, si $\mathcal{X}$ correspond à $\mathcal{M}$, $\mathcal{Y}$ à $\mathcal{N}$, un morphisme $f$ de $\mathcal{X}$ dans $\mathcal{Y}$ sera donné par un foncteur
\[
f^* \;:\; \mathcal{N} \longrightarrow \mathcal{M}
\tag*{(4.1.1)}
\]
en sens inverse, commutant aux $\varinjlim$ inductives quelconques (\footnote{Dire ici que les objets de $\mathcal{M}$ sont visualisés comme des “ouverts” généralisés de $\mathcal{X}$. On a $\mathcal{M}\xrightarrow{\sim}\uHom^!(\mathcal{M}^\circ,\Ens)$, i.e. $\mathcal{M}$ s'interprète bel et bien comme les “faisceaux” sur cette catégorie $\mathcal{M}$ d'ouverts \ldots}). (\textbf{N.B.} si $\mathcal{X}$, $\mathcal{Y}$ sont même des topos, i.e. $\mathcal{M}$, $\mathcal{N}$ des catégories-topos, et si on veut que $f$ soit un morphisme de topos, on exigera de plus que $f^*$ soit de plus \textit{exact à gauche}. On dira qu'un \textit{morphisme de pseudo-topos est exact} si de plus $f^*$ est exact à gauche.)

\msp{page 52}
En vertu du dual de loc. cit. 8.12.7, cela équivaut au fait que $f^*$ admette un adjoint à droite $f_*$,
\[
f_* \;:\; \mathcal{M} \longrightarrow \mathcal{N}\;,
\tag*{(4.1.2)}
\]
lequel commute donc aux petites $\varprojlim$\,. La donnée de $f_*$ satisfaisant à cette propriété détermine donc $f^*$ à isomorphisme unique près. Mais en général, il ne suffit pas de se donner un tel $f_*$, pour pouvoir affirmer qu'il corresponde à un $f^*$, i.e. qu'il admette un adjoint à droite\,; pour ceci, il faudrait qu'on sache que les foncteurs $\mathcal{M}\longrightarrow\Ens$ qui commutent aux petites $\varprojlim$ sont représentables, ce qui sera le cas si $\mathcal{M}$ admet une petite sous-catégorie \textit{co}génératrice, i.e. si $\mathcal{M}^\circ$ est également une catégorie pseudo-topos. % [sic] qu'il admette un adjoint à droite -> adjoint à gauche, f^* is the left adjoint of f_*

\newpage

\printedpage{2}
% head: A. Grothendieck : Dérivateurs, Ch. XVIII, Édition des Universités de Montpellier II et Paris VII
% foot: Version 17 décembre 2025

Il est sans doute judicieux de procéder comme dans le cas des topos, et d'appeler \textit{morphisme de pseudo-topos} $f:\mathcal{X}\longrightarrow\mathcal{Y}$ un quadruple
\[
(f^*,f_*,\alpha,\beta)
\tag*{(4.1.3)}
\]
\msp{page 53}
de deux foncteurs adjoints $(f^*,f_*)$, et des morphismes d'adjonction. Quand $f^*$ est une équivalence, ce qui équivaut à $f_*$ une équivalence, ou $\alpha$ et $\beta$ des isomorphismes, on dira qu'on a une \textit{équivalence de pseudo-topos}, et à toutes fins pratiques, on pourra alors regarder $\mathcal{X}$ et $\mathcal{Y}$ comme des objets essentiellement identiques.

Pour deux pseudo-topos $\mathcal{X}$, $\mathcal{Y}$, j'ai envie d'introduire (comme dans le cas des topos) un pseudo-topos produit
\[
\mathcal{X}\times\mathcal{Y}\;,
\tag*{(4.1.4)}
\]
qui ait les propriétés d'un 2-produit dans la 2-catégorie des pseudo-topos (\footnote{Il apparaît finalement que ce qu'on construit n'est \textit{pas} un 2-produit dans la 2-catégorie envisagée, donc il vaut mieux le noter par un signe tel que $\mathcal{X}\boxtimes\mathcal{Y}$. Mais si $\mathcal{X}$, $\mathcal{Y}$ sont deux topos, $\mathcal{X}\boxtimes\mathcal{Y}$ est un 2-produit dans la 2-catégorie des topos.}). Je devrais préciser que par définition
\[
\begin{aligned}
\uHom_{\pstop}(\mathcal{X},\mathcal{Y}) \;&\simeq\; \uHom_!(\mathcal{F}(\mathcal{Y}),\mathcal{F}(\mathcal{X}))^\circ\\
&\simeq\; \uHom^{!'}(\mathcal{F}(\mathcal{X}),\mathcal{F}(\mathcal{Y})),
\end{aligned}
\tag*{(4.1.5)}
\]
où l'indice ! (resp. l'exposant !) indique la sous-catégorie pleine dans la catégorie $\uHom$ envisagée formée des foncteurs qui commutent aux petites $\varinjlim$ (resp. aux petites
\msp{page 54}
$\varprojlim$\,), et l'accent dans le $\uHom^{!'}$ dans le dernier membre indique qu'on exige de plus que le foncteur en question ait un adjoint à gauche.

Soient donc
\[
\mathcal{M} \;=\; \mathcal{F}(\mathcal{X})\;, \qquad \mathcal{N} \;=\; \mathcal{F}(\mathcal{Y})\;,
\tag*{(4.1.6)}
\]
et soit $\mathcal{Z}$ un troisième pseudo-topos, décrit par
\[
\mathcal{P} \;=\; \mathcal{F}(\mathcal{Z})\;.
\tag*{(4.1.7)}
\]
On veut expliciter, à équivalence près, la catégorie produit $\uHom_{\pstop}(\mathcal{Z},\mathcal{X})\times\uHom_{\pstop}(\mathcal{Z},\mathcal{Y})$, i.e.
\[
\uHom_!(\mathcal{M},\mathcal{P})^\circ \;\times\; \uHom_!(\mathcal{N},\mathcal{Q})^\circ.
\tag*{(4.1.8)}
\] % [sic] \mathcal{Q} -> \mathcal{P}, \mathcal{Q} is not defined; \mathcal{Z} is described by \mathcal{P} (4.1.7)
Je vais introduire pour cela une catégorie
\printedpage{3}
% head: A. Grothendieck : Dérivateurs, Ch. XVIII, Édition des Universités de Montpellier II et Paris VII
% foot: Version 17 décembre 2025
\msp{page 55}
\[
\mathcal{M}\boxtimes\mathcal{N}
\tag*{(4.1.9)}
\]
par la formule
\[
\mathcal{M}\boxtimes\mathcal{N} \;\simeq\; \uHom^{!!}(\mathcal{M}^\circ\times\mathcal{N}^\circ,\Ens)\;,
\tag*{(4.1.10)}
\]
où le signe $\uHom^{!!}$ désigne la sous-catégorie pleine de la catégorie $\uHom$, formée des bi-contrafoncteurs % ?? uncertain: hyphen at line end "bi-|contrafoncteurs" -- bi-contrafoncteurs or bicontrafoncteurs
\[
F \;:\; \mathcal{M}^\circ\times\mathcal{N}^\circ \longrightarrow \Ens\;,
\tag*{(4.1.10.1)}
\]
tels que pour tout $x$ dans $\mathcal{M}$, le foncteur
\[
y \longmapsto F(x,y) \;:\; \mathcal{N}^\circ \longrightarrow \Ens
\tag*{(4.1.10.2)}
\]
soit dans $\uHom^!(\mathcal{N}^\circ,\Ens)$, i.e. transforme petites $\varinjlim$ en $\varprojlim$\,, et que, symétriquement, pour tout $y$ dans $\mathcal{N}$, le foncteur
\[
x \longmapsto F(x,y) \;:\; \mathcal{M}^\circ \longrightarrow \Ens
\tag*{(4.1.10.3)}
\]
soit dans $\uHom^!(\mathcal{M}^\circ,\Ens)$. Mais cela signifie aussi que ces foncteurs sont \textit{représentables}\,; donc on trouve aussi
\[
\mathcal{M}\boxtimes\mathcal{N} \xleftarrow{\sim} \uHom^!(\mathcal{N}^\circ,\mathcal{M})
\tag*{(4.1.11)}
\]
et
\[
\mathcal{M}\boxtimes\mathcal{N} \xleftarrow{\sim} \uHom^!(\mathcal{M}^\circ,\mathcal{N})
\tag*{(4.1.12)}
\]
\msp{page 56}
en utilisant (dans chacune de ces formules) dans deux côtés les équivalences de catégories
\[
\left\{
\begin{aligned}
&\mathcal{M} &&\xrightarrow{\sim} &&\uHom^!(\mathcal{M}^\circ,\Ens)\\
&\mathcal{N} &&\xrightarrow{\sim} &&\uHom^!(\mathcal{N}^\circ,\Ens)\;.
\end{aligned}
\right.
\tag*{(4.1.13)}
\]
On peut dire aussi que la catégorie $\mathcal{M}\boxtimes\mathcal{N}$ est équivalente à celle des \textit{couples} de foncteurs
\[
\left\{
\begin{aligned}
&\varphi &&: &&\mathcal{M}^\circ &&\longrightarrow &&\mathcal{N}\\
&\psi &&: &&\mathcal{N}^\circ &&\longrightarrow &&\mathcal{M}\;,
\end{aligned}
\right.
\tag*{(4.1.14)}
\]
munis de deux morphismes d'adjonction
\[
\left\{
\begin{aligned}
&\psi\varphi^\circ &&\xleftarrow{\alpha} &&\mathrm{id}_{\mathcal{M}}\\
&\varphi\psi^\circ &&\xleftarrow{\beta} &&\mathrm{id}_{\mathcal{N}}
\end{aligned}
\right.
\tag*{(4.1.15)}
\]
satisfaisant les deux compatibilités habituelles exprimant l'adjonction de $\varphi^\circ$ et de $\psi$ (ou, de façon équivalente, de $\psi^\circ$ et de $\varphi$). Le foncteur
\[
(\varphi,\psi) \longmapsto \varphi \qquad (\text{resp.}\;\; (\varphi,\psi) \longmapsto \psi)
\]
donne une équivalence de la catégorie de ces couples avec $\uHom^!(\mathcal{M}^\circ,\mathcal{N})$ (resp. avec $\uHom^!(\mathcal{N}^\circ,\mathcal{M})$), de sorte qu'on a un carré commutatif d'équivalences canoniques de foncteurs :

\end{document}
